% Figure: interaction plot for a two-factor (two-way) experiment.
% Each line is one level of factor B (coating); the x-axis carries the
% levels of factor A (alloy); the y-axis is the cell mean response.
% Parallel lines indicate no interaction; non-parallel lines indicate
% an interaction. Coordinates are the cell means used in the worked
% two-way ANOVA example in the text.
\begin{figure}[htbp]
\centering
\begin{tikzpicture}
\begin{axis}[
    width=11cm, height=6cm,
    xmin=0.7, xmax=3.3,
    ymin=40, ymax=80,
    axis lines=left,
    xlabel={Factor A: alloy},
    ylabel={Cell mean fatigue life (kilocycles)},
    xtick={1,2,3},
    xticklabels={A1, A2, A3},
    ytick={40,50,60,70,80},
    legend style={font=\small, draw=none, fill=none, at={(0.02,0.98)}, anchor=north west},
    every axis plot/.append style={very thick, mark=*, mark size=2pt},
    grid=both,
    grid style={dashed, draw=enggray!20},
]
% Coating B1: roughly parallel-shifted line, modest interaction at A3
\addplot[engblue] coordinates {(1,52) (2,58) (3,66)};
\addlegendentry{coating B1}
% Coating B2: steeper, crosses near A2 -> interaction present
\addplot[engred, dashed] coordinates {(1,48) (2,60) (3,74)};
\addlegendentry{coating B2}
\end{axis}
\end{tikzpicture}
\caption[Interaction plot for a two-factor fatigue experiment]{The
interaction plot reads a two-factor experiment off a single diagram.
Each line connects the cell means for one coating across the three
alloys. Parallel lines would say the coating effect is the same at
every alloy, the additive model. The lines here fan apart toward A3,
the visual signature of an interaction: coating B2 adds little fatigue
life at A1 but a large increment at A3. The formal test is the
interaction $F$-ratio in the two-way analysis of variance worked in
the text; the plot is what the engineer draws first, because an
interaction that the eye sees and the model omits is a
misspecification the residuals will later expose.}
\label{fig:vol02:ch14:twoway-interaction}
\end{figure}
